% ============================================================================
% GLOSSARY.TEX
% Glosarium istilah teknis untuk buku ajar Pemrograman Game
% ============================================================================

\newglossaryentry{game}{
    name={Game},
    description={Aktivitas terstruktur yang melibatkan pemain dengan tujuan, aturan, dan umpan balik}
}

\newglossaryentry{gameobject}{
    name={GameObject},
    description={Objek dasar dalam Unity yang dapat ditempatkan di scene dan memiliki komponen}
}

\newglossaryentry{component}{
    name={Component},
    description={Bagian fungsional yang dapat ditambahkan ke GameObject untuk memberikan perilaku atau properti}
}

\newglossaryentry{prefab}{
    name={Prefab},
    description={Template GameObject yang dapat digunakan kembali dan diinstansiasi berkali-kali}
}

\newglossaryentry{scene}{
    name={Scene},
    description={Konteks atau level dalam game yang berisi GameObject dan pengaturannya}
}

\newglossaryentry{rigidbody}{
    name={Rigidbody},
    description={Komponen Unity yang memungkinkan GameObject dipengaruhi oleh fisika}
}

\newglossaryentry{collider}{
    name={Collider},
    description={Komponen yang mendefinisikan bentuk objek untuk deteksi tumbukan}
}

\newglossaryentry{animator}{
    name={Animator},
    description={Komponen yang mengontrol animasi menggunakan Animator Controller}
}

\newglossaryentry{monobehaviour}{
    name={MonoBehaviour},
    description={Kelas dasar C\# untuk semua skrip Unity yang dapat dilampirkan ke GameObject}
}

\newglossaryentry{coroutine}{
    name={Coroutine},
    description={Fungsi yang dapat menjeda eksekusinya sendiri dan melanjutkan di frame berikutnya}
}

\newglossaryentry{navmesh}{
    name={NavMesh},
    description={Jaringan navigasi yang digunakan untuk pathfinding AI}
}

\newglossaryentry{canvas}{
    name={Canvas},
    description={Area di mana elemen UI Unity dirender}
}

\newglossaryentry{audiosource}{
    name={AudioSource},
    description={Komponen yang memainkan audio clip dalam game}
}

\newglossaryentry{state machine}{
    name={State Machine},
    description={Model komputasi yang terdiri dari state dan transisi antar state}
}

\newglossaryentry{pathfinding}{
    name={Pathfinding},
    description={Proses menemukan jalur optimal dari satu titik ke titik lain}
}

\newglossaryentry{gdd}{
    name={Game Design Document (GDD)},
    description={Dokumen yang menjelaskan konsep, mekanik, dan desain game secara detail}
}

\newglossaryentry{mechanics}{
    name={Mechanics},
    description={Aturan dan sistem yang mengatur bagaimana game berfungsi}
}

\newglossaryentry{dynamics}{
    name={Dynamics},
    description={Perilaku yang muncul dari interaksi mechanics saat dimainkan}
}

\newglossaryentry{aesthetics}{
    name={Aesthetics},
    description={Pengalaman emosional yang dirasakan pemain saat bermain game}
}

\newglossaryentry{instantiate}{
    name={Instantiate},
    description={Membuat instance baru dari GameObject atau Prefab saat runtime}
}

\newglossaryentry{transform}{
    name={Transform},
    description={Komponen yang menyimpan posisi, rotasi, dan skala GameObject}
}

\newglossaryentry{vector3}{
    name={Vector3},
    description={Struktur data yang mewakili titik atau vektor 3D (x, y, z)}
}

\newglossaryentry{quaternion}{
    name={Quaternion},
    description={Representasi matematis rotasi 3D yang menghindari gimbal lock}
}

\newglossaryentry{raycast}{
    name={Raycast},
    description={Teknik untuk mendeteksi objek dengan menembakkan sinar imajiner}
}

\newglossaryentry{object pooling}{
    name={Object Pooling},
    description={Teknik optimasi dengan menggunakan kembali objek daripada membuat dan menghancurkan}
}

\newglossaryentry{lod}{
    name={Level of Detail (LOD)},
    description={Teknik optimasi dengan menggunakan model yang lebih sederhana untuk objek jauh}
}

\newglossaryentry{occlusion culling}{
    name={Occlusion Culling},
    description={Teknik optimasi yang tidak merender objek yang tertutup oleh objek lain}
}

\newglossaryentry{profiler}{
    name={Profiler},
    description={Alat Unity untuk menganalisis performa game}
}

\newglossaryentry{delegate}{
    name={Delegate},
    description={Tipe yang merepresentasikan referensi ke method dengan signature tertentu}
}

\newglossaryentry{event}{
    name={Event},
    description={Mekanisme untuk komunikasi antar objek tanpa coupling langsung}
}

\newglossaryentry{async}{
    name={Async/Await},
    description={Pola pemrograman asinkron untuk operasi non-blocking}
}

\newglossaryentry{blend tree}{
    name={Blend Tree},
    description={Struktur dalam Animator Controller untuk transisi halus antar animasi}
}

\newglossaryentry{hud}{
    name={HUD (Heads-Up Display)},
    description={Elemen UI yang ditampilkan di atas gameplay (health, score, dll)}
}

\newglossaryentry{ui}{
    name={UI (User Interface)},
    description={Elemen antarmuka yang memungkinkan interaksi pemain dengan game}
}

\newglossaryentry{ux}{
    name={UX (User Experience)},
    description={Pengalaman keseluruhan pemain saat berinteraksi dengan game}
}

\newglossaryentry{sfx}{
    name={SFX (Sound Effects)},
    description={Efek suara yang dipicu oleh aksi dalam game}
}

\newglossaryentry{bgm}{
    name={BGM (Background Music)},
    description={Musik latar yang diputar selama gameplay}
}

\newglossaryentry{audio mixer}{
    name={Audio Mixer},
    description={Sistem Unity untuk mengontrol dan memproses audio}
}

\newglossaryentry{spawn}{
    name={Spawn},
    description={Proses membuat objek baru dalam game pada waktu tertentu}
}

\newglossaryentry{collectible}{
    name={Collectible},
    description={Item dalam game yang dapat dikumpulkan oleh pemain}
}

\newglossaryentry{checkpoint}{
    name={Checkpoint},
    description={Titik dalam game di mana pemain dapat melanjutkan jika gagal}
}

\newglossaryentry{npc}{
    name={NPC (Non-Player Character)},
    description={Karakter dalam game yang dikendalikan oleh AI, bukan pemain}
}

\newglossaryentry{ai}{
    name={AI (Artificial Intelligence)},
    description={Kecerdasan buatan yang mengendalikan perilaku karakter atau sistem dalam game}
}

\newglossaryentry{platform}{
    name={Platform},
    description={Sistem operasi atau perangkat tempat game dijalankan (PC, Mobile, Web)}
}

\newglossaryentry{build}{
    name={Build},
    description={Proses mengkompilasi project Unity menjadi executable untuk platform tertentu}
}

\newglossaryentry{deployment}{
    name={Deployment},
    description={Proses mendistribusikan game yang sudah di-build ke platform distribusi}
}

\newglossaryentry{version control}{
    name={Version Control},
    description={Sistem untuk melacak perubahan kode dan kolaborasi tim}
}

\newglossaryentry{git}{
    name={Git},
    description={Sistem version control terdistribusi yang populer}
}

\newglossaryentry{api}{
    name={API (Application Programming Interface)},
    description={Antarmuka yang memungkinkan komunikasi antara komponen software}
}

\newglossaryentry{unity hub}{
    name={Unity Hub},
    description={Aplikasi manajemen untuk menginstal dan mengelola versi Unity yang berbeda}
}

\newglossaryentry{asset store}{
    name={Asset Store},
    description={Marketplace Unity untuk membeli dan mengunduh aset, tools, dan ekstensi}
}

\newglossaryentry{package}{
    name={Package},
    description={Koleksi file dan resource yang dapat diimpor ke project Unity}
}

\newglossaryentry{script}{
    name={Script},
    description={File kode C\# yang berisi logika untuk mengontrol perilaku GameObject}
}

\newglossaryentry{runtime}{
    name={Runtime},
    description={Waktu ketika game sedang berjalan, bukan saat editing}
}

\newglossaryentry{editor}{
    name={Editor},
    description={Lingkungan pengembangan Unity untuk membuat dan mengedit game}
}

\newglossaryentry{inspector}{
    name={Inspector},
    description={Panel Unity yang menampilkan properti dan komponen GameObject yang dipilih}
}

\newglossaryentry{hierarchy}{
    name={Hierarchy},
    description={Panel Unity yang menampilkan struktur GameObject dalam scene}
}

\newglossaryentry{project}{
    name={Project},
    description={Panel Unity yang menampilkan struktur file dan folder project}
}

\newglossaryentry{console}{
    name={Console},
    description={Panel Unity yang menampilkan pesan log, warning, dan error}
}
